\begin{frame}{자주 하는 실수: 롤아웃 정책과 시뮬레이션 횟수}
  \begin{itemize}
    \item \textbf{① 롤아웃 정책을 너무 정교하게 만들려는 것}
      --- ``무작위 롤아웃은 부정확하니 정교한 규칙을 넣자''는
      생각에 복잡한 로직을 넣으면, 시뮬레이션 \textbf{한 번의
      계산 비용}이 커져 같은 시간 안에 \textbf{훨씬 적은
      횟수}만 돎 --- 위 님 실습처럼 \textbf{``단순하지만 많이''
      도는 쪽이} 나을 때가 많음
      \begin{itemize}
        \item AlphaZero가 신경망으로 롤아웃을 대체한 것도
          ``더 똑똑하게''가 아니라 ``\textbf{한 번의 평가로
          승패 시뮬레이션 수백 번을 대체}할 만큼 신뢰도 높은
          신호를 준다''는 \textbf{효율성 계산}
      \end{itemize}
    \item \textbf{② UCT 탐험 상수 $c$를 문제와 무관하게
      고정}하는 것 --- $c=1.4$는 님 게임에는 잘 맞았지만,
      승/패 대신 \textbf{연속적인 점수}를 다루는 문제에서는
      값의 \textbf{스케일}이 다르므로 $c$를 다시 튜닝해야
      (Week 2.2에서 UCB의 $c$를 데이터 스케일에 맞춰
      조정해야 했던 것과 같은 이유)
  \end{itemize}
\end{frame}
